%قبل از شروع مسئله ترجیحاً چیزی ننویسید. اما اگر خواستید بنویسید، در قالب لیستی (فقرات و شمارش) بنویسید.
\شروع‌مسئله[12]  
% صورت اصلی مسئله را در این‌جا بنویسید.
توابع هر دسته را بر اساس پیچیدگی‌شان مرتب کنید. (فقط ترتیب را از کوچک به بزرگ بنویسید؛ نیاز به نوشتن اثبات نیست.)
\شروع{زیرمسائل}
\زیرمسئله[4]
\begin{large}
		\begin{center}
			\begin{tabular}{l l}
				$f_2(n) = 100! $&
				$f_1(n) = \binom{n}{n-3} $\\
				$f_4(n) = n\sqrt[3]{n}$&
				$f_3(n) = \pi^n $
			\end{tabular}
		\end{center}
\end{large}

\زیرمسئله[4]
\begin{large}
		\begin{center}
			\begin{tabular}{l l}
				$f_2(n) = n^5\cdot5^n $&
				$f_1(n) = n\sum\limits_{i=1}^n \frac{1}{\sqrt{i} + \sqrt{i+1}} $\\
				$f_4(n) = \sqrt{2^{\sqrt{n}}}$&
				$f_3(n) = 2^{\sqrt[n]{n^{n+1}}} $
			\end{tabular}
		\end{center}
\end{large}

\زیرمسئله[4]

\begin{large}
		\begin{center}
			\begin{tabular}{l l}
				$f_2(n) = n\sqrt{n} $&
				$f_1(n) = 2^{\log_2 n^2} $\\
				$f_4(n) = \sqrt{n}\log n$&
				$f_3(n) = n\binom{n}{2} $
			\end{tabular}
		\end{center}
\end{large}

\پایان{زیرمسائل}
%اگر مسئله چند قسمت داشت به این صورت عمل کنید:
%\شروع{زیرمسائل}
%\زیرمسئله[0] % مسئلۀ ۱آ
%\زیرمسئله[0] % مسئلۀ ۱ب
%\زیرمسئله[0] % مسئلۀ ۱پ
%\زیرمسئله[0] % مسئلۀ ۱ت
%\پایان{زیرمسائل}

\solution{
آ. 
$f_2 < f_4 < f_1 < f_3$

ب.
$f_1 < f_4 < f_2 < f_3$

پ.
$f_4 < f_2 < f_1 < f_3$

}




%%%%%%%%%%%%%%%%%%%%%%%%%%%%%%%%%%%%%%%%%%%%%%%%%%%%%
% برخی دستورات مفید برای تمرین‌ها. %
%%%%%%%%%%%%%%%%%%%%%%%%%%%%%%%%%%%%%%%%%%%%%%%%%%%%%

% فرمول‌های نمونه:
%$f(x) = \Theta(x)$
%$T(x, y) \leq \log(x) + 2^y + \binom{2n}{n}$
% {\tt code\_function}


% لیست موارد بدون شماره‌گذاری:
%\شروع{فقرات}
%    \فقره مورد اول 
%        \شروع{فقرات}
%            \فقره مورد اول داخلی 
%                \شروع{فقرات}
%                    \فقره اولین مورد خیلی داخلی 
%                    \فقره دومین مورد خیلی داخلی 
%                \پایان{فقرات}
%            \فقره مورد دوم داخلی 
%        \پایان{فقرات}
%    \فقره مورد دوم 
%\پایان{فقرات}

% لیست موارد شماره‌گذاری شده:
%\شروع{شمارش}
%    \فقره یک مورد 
%    \فقره یک مورد دیگر 
%    \فقره باز هم یک مورد دیگر
%\پایان{شمارش}

% برای نوشتن یک فرمول چند خطی که یک شماره بخوره. برای حالت‌های دیگر به سایت زیر مراجعه کنید:
% https://www.overleaf.com/learn/latex/Aligning_equations_with_amsmath
%\begin{align} 
%    \begin{split}
%        (x+y)^3 &= (x+y)^2(x+y) \\
%                &= (x^2+2xy+y^2)(x+y) \\
%                &= (x^3+2x^2y+xy^2) + (x^2y+2xy^2+y^3) \\
%                &= x^3+3x^2y+3xy^2+y^3
%    \end{split}                                 
%\end{align}

% نوشتن ماتریس و آرایه دوبعدی
%\begin{align}
%    A = 
%    \begin{bmatrix}
%        A_{11} & A_{21} \\
%        A_{21} & A_{22}
%    \end{bmatrix}
%\end{align}

% اضافه کردن تصویر با شرح. برای اضافه کردن تصویر بدون شرح، فقط دستور درج‌تصویر را استفاده کنید.
%\شروع{شکل}[h!]
%\تنظیم‌ازوسط
%  \درج‌تصویر[پهنا=0.5\پهنای‌متن]{img.jpg}
%\شرح{شرح تصویر}
%\برچسب{تصویر:۱}
%\پایان{شکل}
%\پایان{زیرمسائل}

%کد زیر را ببینید.
%%\lstinputlisting{source_filename.py}
%\begin{latin}
%\begin{lstlisting}
%# Import the modules
%import sys
%import random
%
%ans = True
%
%while ans:
%question = raw_input("Ask the magic 8 ball a question: (press enter to quit) ")
%
%answers = random.randint(1,8)
%
%if question == "":
%sys.exit()
%
%elif answers == 1:
%print "It is certain"
%
%elif answers == 2:
%print "Outlook good"
%
%elif answers == 3:
%print "You may rely on it"
%
%elif answers == 4:
%print "Ask again later"
%
%elif answers == 5:
%print "Concentrate and ask again"
%
%elif answers == 6:
%print "Reply hazy, try again"
%
%elif answers == 7:
%print "My reply is no"
%
%elif answers == 8:
%print "My sources say no"
%\end{lstlisting}
%\end{latin}
